% GENERATED FILE. Edit canonical lessons, then run npm run generate:books.
% canonical-source-hash: fnv1a64:687d733dbcca55bd
% canonical-lessons: GE-C440-buergersteig, GE-C440-parkbank, GE-C440-laterne, GE-C440-kopfsteinpflaster, GE-C440-brunnen

\chapter{Auf der Straße}
\label{ch:ge-auf-der-stra-e}
\hlchaptermodality{440}

\begin{chapteropening}
\textbf{By the end of this chapter:} \emph{I can describe things in a street: the pavement, the park bench, the street lamp, the cobbles and the fountain.}

Complete the last lesson of chapter 440: I can describe things in a street: the pavement, the park bench, the street lamp, the cobbles and the fountain.
\end{chapteropening}

\section[der Bürgersteig]{der Bürgersteig — the pavement (US: sidewalk)}

\label{lesson:GE-C440-buergersteig}



\begin{quote}
Before the new one: say the German for the police station, then the German for the barracks.
\end{quote}

\subsection*{You'll want to know: der Bürgersteig}
\textbf{der Bürgersteig} — \textquotedblleft{}the pavement (US: sidewalk)\textquotedblright{}.

\begin{quote}
Die Kinder fahren mit dem Roller auf dem Bürgersteig.
\end{quote}

\begin{quote}
\emph{The children ride their scooters on the pavement.}
\end{quote}

Plural: \emph{die Bürgersteige}. Also \emph{der Gehweg}; in Switzerland \emph{das Trottoir}.

\begin{cousinweb}
\emph{Bürgersteig} is built from \emph{Bürger} 'townsman, citizen', from \emph{Burg} 'castle, town', and \emph{Steig} 'path', inherited from Old High German \emph{stīg} (EtymDB).
\end{cousinweb}

\begin{grammarlens}[title={What you've built}]
One more word for telling, explaining and describing at B1.
\end{grammarlens}

\subsection*{Your turn}
\begin{itemize}
\raggedright
  \item \emph{Say it:} \emph{der Bürgersteig}
  \item \emph{Say it:} \emph{der Bürgersteig}, once more
  \item \emph{Say it:} a sentence of your own with \emph{der Bürgersteig}
  \item \emph{Recall:} say the German for the police station, then the German for the barracks, then say \emph{der Bürgersteig} again
\end{itemize}

\begin{tcolorbox}[breakable,colback=teal!4,colframe=teal!35!black,title={Before you move on}]
What does der Bürgersteig mean? (\textquotedblleft{}The pavement (US: sidewalk)\textquotedblright{}.) Say it once more.
\end{tcolorbox}
\section[die Parkbank]{die Parkbank — the park bench}

\label{lesson:GE-C440-parkbank}



\begin{quote}
Before the new one: say the German for the barracks, then the German for the pavement.
\end{quote}

\subsection*{You'll want to know: die Parkbank}
\textbf{die Parkbank} — \textquotedblleft{}the park bench\textquotedblright{}.

\begin{quote}
Ein alter Mann füttert auf der Parkbank die Tauben.
\end{quote}

\begin{quote}
\emph{An old man feeds the pigeons on the park bench.}
\end{quote}

Plural: \emph{die Parkbänke}. \emph{Die Bank} is a bench (plural \emph{Bänke}) or a bank (plural \emph{Banken}).

\begin{cousinweb}
\emph{Parkbank} is built from \emph{Park}, which EtymDB traces to Old French \emph{parc} 'enclosure, pen', and \emph{Bank} 'bench', inherited from Old High German \emph{bank} (EtymDB). The money \emph{Bank} came back into German from Italian \emph{banca}, a money-changer's counter.
\end{cousinweb}

\begin{grammarlens}[title={What you've built}]
One more word for telling, explaining and describing at B1.
\end{grammarlens}

\subsection*{Your turn}
\begin{itemize}
\raggedright
  \item \emph{Say it:} \emph{die Parkbank}
  \item \emph{Say it:} \emph{die Parkbank}, once more
  \item \emph{Say it:} a sentence of your own with \emph{die Parkbank}
  \item \emph{Recall:} say the German for the barracks, then the German for the pavement, then say \emph{die Parkbank} again
\end{itemize}

\begin{tcolorbox}[breakable,colback=teal!4,colframe=teal!35!black,title={Before you move on}]
What does die Parkbank mean? (\textquotedblleft{}The park bench\textquotedblright{}.) Say it once more.
\end{tcolorbox}
\section[die Laterne]{die Laterne — the street lamp; the lantern}

\label{lesson:GE-C440-laterne}



\begin{quote}
Before the new one: say the German for the pavement, then the German for the park bench.
\end{quote}

\subsection*{You'll want to know: die Laterne}
\textbf{die Laterne} — \textquotedblleft{}the street lamp; the lantern\textquotedblright{}.

\begin{quote}
Die Laternen gehen um sechs Uhr an.
\end{quote}

\begin{quote}
\emph{The street lamps come on at six o'clock.}
\end{quote}

Plural: \emph{die Laternen}. Also a paper lantern.

\begin{cousinweb}
\emph{Laterne} was borrowed from Latin \emph{lanterna} 'lamp, lantern', also spelled \emph{lāterna} (Lewis \& Short). EtymDB lists the German word without a source.
\end{cousinweb}

\begin{grammarlens}[title={What you've built}]
One more word for telling, explaining and describing at B1.
\end{grammarlens}

\subsection*{Your turn}
\begin{itemize}
\raggedright
  \item \emph{Say it:} \emph{die Laterne}
  \item \emph{Say it:} \emph{die Laterne}, once more
  \item \emph{Say it:} a sentence of your own with \emph{die Laterne}
  \item \emph{Recall:} say the German for the pavement, then the German for the park bench, then say \emph{die Laterne} again
\end{itemize}

\begin{tcolorbox}[breakable,colback=teal!4,colframe=teal!35!black,title={Before you move on}]
What does die Laterne mean? (\textquotedblleft{}The street lamp; the lantern\textquotedblright{}.) Say it once more.
\end{tcolorbox}
\section[das Kopfsteinpflaster]{das Kopfsteinpflaster — the cobbles, the cobblestone paving}

\label{lesson:GE-C440-kopfsteinpflaster}



\begin{quote}
Before the new one: say the German for the park bench, then the German for the street lamp.
\end{quote}

\subsection*{You'll want to know: das Kopfsteinpflaster}
\textbf{das Kopfsteinpflaster} — \textquotedblleft{}the cobbles, the cobblestone paving\textquotedblright{}.

\begin{quote}
Auf dem Kopfsteinpflaster ist das Fahrrad sehr laut.
\end{quote}

\begin{quote}
\emph{The bike is very noisy on the cobbles.}
\end{quote}

No plural. \emph{Der Kopfstein} is one stone, round like a head.

\begin{cousinweb}
\emph{Kopfsteinpflaster} is built from \emph{Kopf} 'head', \emph{Stein} 'stone', inherited from Proto-Germanic \emph{stainaz} (EtymDB), and \emph{Pflaster} 'paving', which goes back to Latin \emph{emplastrum} 'plaster' (EtymDB).
\end{cousinweb}

\begin{grammarlens}[title={What you've built}]
One more word for telling, explaining and describing at B1.
\end{grammarlens}

\subsection*{Your turn}
\begin{itemize}
\raggedright
  \item \emph{Say it:} \emph{das Kopfsteinpflaster}
  \item \emph{Say it:} \emph{das Kopfsteinpflaster}, once more
  \item \emph{Say it:} a sentence of your own with \emph{das Kopfsteinpflaster}
  \item \emph{Recall:} say the German for the park bench, then the German for the street lamp, then say \emph{das Kopfsteinpflaster} again
\end{itemize}

\begin{tcolorbox}[breakable,colback=teal!4,colframe=teal!35!black,title={Before you move on}]
What does das Kopfsteinpflaster mean? (\textquotedblleft{}The cobbles, the cobblestone paving\textquotedblright{}.) Say it once more.
\end{tcolorbox}
\section[der Brunnen]{der Brunnen — the fountain; the well}

\label{lesson:GE-C440-brunnen}



\begin{quote}
Before the new one: say the German for the street lamp, then the German for the cobbles.
\end{quote}

\subsection*{You'll want to know: der Brunnen}
\textbf{der Brunnen} — \textquotedblleft{}the fountain; the well\textquotedblright{}.

\begin{quote}
Auf dem Marktplatz steht ein alter Brunnen.
\end{quote}

\begin{quote}
\emph{There is an old fountain in the market square.}
\end{quote}

Same form in the plural. \emph{Der Springbrunnen} — a fountain that sprays.

\begin{cousinweb}
\emph{Brunnen} is inherited through Middle High German \emph{brunne} from Old High German \emph{brunno} 'spring, well' and Proto-Germanic \emph{brunnô} (EtymDB). Scots \emph{burn} 'stream' is a cousin.
\end{cousinweb}

\begin{grammarlens}[title={What you've built}]
One more word for telling, explaining and describing at B1.
\end{grammarlens}

\subsection*{Your turn}
\begin{itemize}
\raggedright
  \item \emph{Say it:} \emph{der Brunnen}
  \item \emph{Say it:} \emph{der Brunnen}, once more
  \item \emph{Say it:} a sentence of your own with \emph{der Brunnen}
  \item \emph{Recall:} say the German for the street lamp, then the German for the cobbles, then say \emph{der Brunnen} again
\end{itemize}

\begin{tcolorbox}[breakable,colback=teal!4,colframe=teal!35!black,title={Before you move on}]
What does der Brunnen mean? (\textquotedblleft{}The fountain; the well\textquotedblright{}.) Say it once more.
\end{tcolorbox}
